\section{Background}

The concept of \textit{Recursive Spawning} emerges at the intersection of four distinct but converging research threads: agent lifecycle management in multi-agent systems, accountability and provenance infrastructure for distributed AI, the design tradeoffs between fixed and mutable delegation chains, and hierarchical task decomposition as a planning paradigm. Each thread has developed independently, and each carries lessons relevant to the structural problem Recursive Spawning addresses. This section surveys the relevant landscape and positions the BX3 Framework as a substrate capable of unifying these threads under a single architectural logic.

% ----------------------------------------------------------------
\subsection{Agent Lifecycle Management in Multi-Agent Systems}

Modern AI agent ecosystems have evolved from single-agent runners to orchestrated populations in which agents are created, configured, delegated to, and retired over time. Managing this lifecycle rigorously is an active area of research and standards development.

The \textbf{Agent Lifecycle Protocol (ALP)} \cite{moore2025hmas} defines six canonical lifecycle events---Genesis, Fork, Migration, Retraining, Succession, and Decommission---with explicit state machine semantics, transition rules, and boundary hooks. ALP's central contribution is the \textit{Fork Registry}, which maintains a lineage history for every agent, including genetic provenance (the agent's origin) and epigenetic provenance (the accumulated experience and configuration changes that shape its current state). This registry enables any party to reconstruct an agent's full family tree, configurations, and reputation history at any point in time.

The \textbf{Agent Identity and Trust Lifecycle Protocol (AITLP)} \cite{trusttrack2025} adds governance-oriented lifecycle concepts, emphasizing auditable actions, drift detection, and bounded mandates. AITLP introduces the concept of \textit{Agent Legacy Mode} for generational knowledge transfer across agent succession events, ensuring that a retiring agent's institutional knowledge is preserved and transferable to its successor.

Within production-grade agent architectures, the emphasis is on fail-safe behavior, sandbox-first execution, deterministic fallback workflows, and human approval gates for risky actions \cite{alenezi2026}. BX3 incorporates these concerns as named architectural pillars---most notably the Sandbox Gate and Bailout Protocol---rather than as ad hoc implementation heuristics. The \textbf{Tiered Agentic Oversight (TAO)} framework \cite{kim2025tao} demonstrates that hierarchical supervision among specialized agents can reduce error propagation in safety-critical settings; BX3 distinguishes itself by making human accountability the unconditional terminal endpoint of unresolved escalation rather than a high-tier supervisory option.

The key insight from lifecycle management literature is that lifecycle transitions are not merely operational events---they are governance moments at which accountability must be explicitly transferred. Whether an agent is forked, migrated, or decommissioned, the system must record who held accountability before the transition and who holds it after. This requirement is architectural, not optional.

% ----------------------------------------------------------------
\subsection{Accountability and Provenance in Distributed Systems}

Accountability in distributed AI systems requires more than logging. It demands a traceable, tamper-evident record of decision authority, execution provenance, and outcome attribution that survives the complexity of multi-agent choreography.

The \textbf{Human Delegation Provenance (HDP) Protocol} \cite{hdp2025arxiv} addresses this with a lightweight, token-based cryptographic mechanism that creates an append-only chain of signed hops from human authorization through every agent in a delegation path. The key property of HDP is \textit{offline verifiability}: any participant can validate the entire provenance chain using only the issuer's public key and a session identifier, without relying on external registries or trust anchors. HDP's model is close in spirit to what BX3 requires, but HDP operates at the protocol level without prescribing the internal structure of the agents or layers it connects.

The \textbf{Agent Execution Record (AER)} \cite{proml2022} introduces a first-class structured provenance primitive that captures intent, observations, inferences, plans with revision rationale, evidence chains, and confidence-scored verdicts at every step of an agent's investigation. Crucially, AER records \textit{delegation authority chains}---who delegated what to whom, and how decisions propagated across the chain---making the delegation graph itself queryable. This delegation authority chain is the most directly relevant predecessor to BX3's Immutable Parent Pointer, though AER does not enforce the unconditional escalation property that BX3 requires.

The \textbf{SentinelAgent} framework \cite{sentinelagent2026} formalizes the problem through the \textit{Delegation Chain Calculus (DCC)}, which defines seven properties of secure delegation chains. Six of these are deterministic: \textbf{authority narrowing} (delegated authority cannot exceed the delegator's), \textbf{policy preservation} (policy constraints propagate unchanged), \textbf{forensic reconstructibility} (any action can be traced back through the chain), \textbf{cascade containment} (delegation cannot exceed its declared scope), \textbf{scope-action conformance} (actions stay within scope), and \textbf{output schema conformance} (outputs match declared formats). The seventh property, \textbf{intent preservation}, is probabilistic and represents the hard problem of ensuring a delegate acts in the spirit of the delegation, not merely its letter. SentinelAgent's emphasis on forensic reconstructibility and cascade containment is directly aligned with the upstream accountability guarantee that BX3 encodes in its Immutable Parent Pointer. However, SentinelAgent treats the human anchor as one node among many in the delegation chain; BX3's contribution is to make the human anchor architecturally required and unreachable by any other path when a child node exhausts its bounds.

The broader provenance literature is converging on the same requirement. The \textbf{PROV-AGENT} framework \cite{provcagent2025} extends the W3C PROV standard with the Model Context Protocol to capture fine-grained, agent-centric metadata (prompts, responses, decisions) linked to end-to-end workflow context. The DZone accountability synthesis \cite{dzone2025accountability} summarizes the consensus: when an agent's action leads to a bad outcome, there must be a traceable, identifiable accountable party who can explain the decision, and regulators and standards bodies (NIST AI RMF, EU AI Act) increasingly mandate governance, traceability, auditability, documented roles, risk ownership, decision provenance, logging, and human oversight for AI-driven systems.

% ----------------------------------------------------------------
\subsection{Fixed vs. Mutable Delegation Chains}

A fundamental design decision in any delegation architecture is whether the delegation chain itself is fixed at creation or mutable during execution. This choice has profound consequences for accountability, auditability, and the system's ability to recover from failures.

\textbf{Fixed delegation chains} establish the complete delegation path at spawn time and do not allow it to change during execution. Each child node receives a reference to its parent, and that reference is immutable. This approach has the advantage of making provenance trivial to reconstruct: the entire chain is known a priori and can be audited in full. HDP's append-only signed hop model is a form of fixed delegation, in which each hop is recorded but the path is never altered after the fact. Fixed chains are predictable and enforceable; they also simplify the application of policies like authority narrowing and scope-action conformance, because the complete boundary of the delegation is known at the moment the chain is formed.

However, fixed delegation chains carry a significant liability in open-ended or long-running tasks: they cannot adapt. If a delegated subtask fails, the parent cannot re-delegate to a different agent without breaking the immutability of the chain. If a more appropriate specialist becomes available mid-execution, the fixed chain cannot incorporate them. These limitations have motivated a body of work on \textbf{mutable delegation chains} that allow the delegation graph to evolve during execution.

The \textbf{IACT framework} \cite{agenthive2025} advocates a dynamic, runtime-adaptive agent topology that grows incrementally to fit a problem, replacing fixed chains of function calls with bidirectional, stateful dialogues that enable ongoing verification, error correction, and ambiguity resolution. Under IACT, agent lineages (parent-child relationships) evolve during execution rather than being predetermined by developers. The benefit is mitigation of error propagation common in linear, unidirectional chains; the cost is that provenance becomes order-dependent and harder to audit post-hoc, because the final delegation graph is only knowable after execution completes.

The \textbf{Layered Mutability} analysis \cite{layeredmutability2026} identifies five layers of agent change---pretraining, post-training alignment, self-narrative, memory, and weight-level adaptation---and shows that reversion at one layer (e.g., resetting a self-description after memory accumulation) does not restore baseline behavior. The system exhibits \textit{identity hysteresis}: accumulated drift persists even after attempted rollback. This finding has direct implications for mutable delegation chains: if a chain can be rewritten mid-execution, the rewritten chain is not the only record of what happened. The original chain must be preserved as an audit artifact, which effectively requires the mutability to be implemented as an append-only log of mutations rather than true in-place replacement.

The architecture that emerges from this tension is a \textbf{structurally immutable, operationally flexible} delegation chain: the parent pointer structure is immutable and permanent, but the tasks and context flowing through that structure can be dynamically reallocated, monitored, and recovered. This is the design principle that BX3 encodes in its Immutable Parent Pointer. The chain of accountability is fixed and auditable; the work within the chain is not.

% ----------------------------------------------------------------
\subsection{Hierarchical Task Decomposition in AI}

Recursive Spawning is fundamentally a hierarchical task decomposition mechanism: a parent node decomposes a task into sub-tasks, spawns child nodes to execute each sub-task, and synthesizes the results into a coherent whole. This pattern is well-established in the agentic AI literature, with variations that differ primarily in how hierarchical the decomposition is, how dynamic the spawning is, and how the parent retains oversight.

\textbf{ReAcTree} \cite{reactree2025} presents a hierarchical framework for long-horizon task planning that decomposes complex goals into subgoals using a dynamic agent tree. Each subgoal is handled by an LLM agent node capable of reasoning, acting, and further expanding the tree, while control-flow nodes coordinate overall execution. ReAcTree introduces two memory systems: episodic-memory-style goal-specific subgoal examples for each agent node and working-memory-style environment observations shared across nodes. The key architectural feature is that parent nodes delegate and spawn child agent nodes, with control-flow nodes managing execution strategies across the tree.

The \textbf{Agent Contracts} framework \cite{agentcontracts2025} formalizes the delegation relationship as a bounded optimization problem: an agent receives a contract specifying resource budgets (tokens, API calls), time limits (deadlines, durations), and output specifications. Contracts are applied to parent-child relationships, where each child agent receives a bounded scope and budget, ensuring controlled delegation and preventing runaway tasks. The immutability of contract terms is the mechanism by which delegation chains achieve reproducibility and auditability; each parent can enforce compliance by distributing constrained work packages to children and aggregating results under a unified contract.

The \textbf{FIRMHIVE} framework \cite{firmhive2025} transforms delegation into executable primitives, building a runtime \textit{Tree of Agents (ToA)} for coordinated reasoning. Each agent performs per-agent, executable tasks, enabling dynamic parent-child delegation chains within the hierarchy rather than static, hand-crafted workflows. FIRMHIVE's empirical results demonstrate that hierarchical, dynamic delegation achieves deeper and broader analysis than static workflows: roughly 16 times more reasoning steps, 2.3 times more files inspected, and 5.6 times more alerts per firmware image compared to baseline LLM-agent systems.

At the most general level, the \textbf{adaptive delegation} literature \cite{agenthive2025} argues for moving beyond simple heuristics to dynamically adapt task allocation and sub-task decomposition in complex, changing environments. This work emphasizes explicit handling of authority, responsibility, accountability, and trust within delegation networks; clear role definitions, boundaries, and transfer of control; and continuous monitoring and feedback-driven adjustments to ensure task completion under specified constraints. The framework supports both centralized and decentralized orchestration, addressing safety and reliability in high-stakes deployments.

What unites these approaches is that hierarchical decomposition is treated as a first-class architectural concern, not an implementation detail. The decomposition strategy, the delegation semantics, and the result synthesis are all explicit and auditable. BX3 inherits this commitment and adds the structural constraint that the decomposition cannot create orphaned nodes: every child must have a traceable parent, and every parent chain must eventually reach a human anchor.

% ----------------------------------------------------------------
\subsection{The BX3 Framework as Substrate for Immutable Parent Chains}

The background reviewed above establishes that the problem of accountable delegation in multi-agent AI systems is well-identified, actively studied, and architecturally maturing. What remains unsolved is the structural constraint that makes accountability unconditionally upstream: no system reviewed here guarantees that every spawned node has a traceable path to a human accountable party, and that every unresolved escalation terminates at that party rather than propagating further into the machine stack.

The BX3 Framework provides the missing structural layer. Its three-layer architecture---Purpose Layer (accountability), Bounds Engine (bounded reasoning), and Fact Layer (deterministic enforcement)---is not merely a categorization of agent capabilities but an \textit{immutability architecture}. Each layer has fixed functional properties that cannot be varied without violating the layer's identity. When a child node is spawned under BX3, the parent pointer is established not as a policy choice but as a structural necessity: the child exists within a layered architecture, and the topmost layer of that architecture is always a human anchor. The Immutable Parent Pointer makes this guarantee physically enforceable rather than policy-dependent.

The Recursive Spawning mechanism, which Section 3 defines formally, builds on this substrate. It treats spawning as an operation that must satisfy the layer's immutability constraints at the moment of creation. The parent pointer is not an optional metadata field---it is a structural binding that propagates the full BX3 stack to every child, ensuring that each spawned node has a Purpose, a Bounds Engine, and a Fact Layer, and that unresolved escalations move upward through the hierarchy until they reach the human anchor. This is not a feature added to existing agent architectures; it is the architectural foundation on which those architectures must be built to make accountability structural rather than optional.